\documentclass[tikz,border=10pt]{standalone}
\usetikzlibrary{positioning, arrows.meta, calc, patterns, backgrounds, babel}
\usepackage[utf8]{inputenc}
\usepackage[T1]{fontenc}
\usepackage{lmodern}
\usepackage[spanish,es-tabla]{babel}
\usepackage{amsmath, amssymb, bm}

\begin{document}
\begin{tikzpicture}[
  align=center,
  font=\sffamily\small,
  box/.style={
    draw=black!70, rectangle, rounded corners=3pt,
    minimum width=3.2cm, minimum height=1.2cm,
    thick, fill=white, inner sep=6pt
  },
  widebox/.style={box, minimum width=3.8cm},
  arr/.style ={draw, -{Stealth[scale=1.15]}, thick},
  darr/.style={draw, -{Stealth[scale=1.15]}, thick, dashed, draw=blue!60},
  tarr/.style={draw, -{Stealth[scale=1.05]}, semithick, dashed, draw=black!45},
]

%% ====================================================================
%%  GEOMETRIA  (coordenadas absolutas ensanchadas)
%%
%%  Bandas:  x in [-9.0, +8.5]  (ancho = 17.5 cm)
%%
%%  y_ctrl =  0      semialtura = 1.90
%%  y_fis  = -4.60   semialtura = 1.50
%%  y_tel  = -7.90   semialtura = 0.90
%%
%%  Columnas de bloques:
%%    ref/obs : x = -6.0
%%    ctrl    : x = -1.0
%%    mix     : x =  5.2
%%    act     : x =  5.2   (mismo x que mix, pero en banda fisica)
%%    rk4     : x = -2.2
%%    tel     : x =  2.0   (desplazado a la derecha para dejar paso a la izq)
%%
%%  Canales exteriores:
%%    canal ZOH     x =  8.0
%%    canal FB      x = -8.4
%% ====================================================================

\def\xBL{-9.0}   \def\xBR{8.5}
\def\yC{0}     \def\hC{1.90}
\def\yF{-4.60} \def\hF{1.50}
\def\yT{-7.90} \def\hT{0.90}
\def\xZOH{9.0}
\def\xFB{-8.4}

%% ====================================================================
%%  BANDAS DE FONDO
%% ====================================================================
\begin{scope}[on background layer]
  \filldraw[draw=blue!45, fill=blue!3,
            pattern=vertical lines, pattern color=blue!8,
            rounded corners=6pt, dashed, line width=1pt]
    (\xBL,\yC-\hC) rectangle (\xBR,\yC+\hC);

  \filldraw[draw=green!50, fill=green!3,
            pattern=horizontal lines, pattern color=green!9,
            rounded corners=6pt, dashed, line width=1pt]
    (\xBL,\yF-\hF) rectangle (\xBR,\yF+\hF);

  \filldraw[draw=orange!50, fill=orange!3,
            pattern=north west lines, pattern color=orange!8,
            rounded corners=6pt, dashed, line width=1pt]
    (\xBL,\yT-\hT) rectangle (\xBR,\yT+\hT);
\end{scope}

%% ====================================================================
%%  ETIQUETAS DE BANDA
%% ====================================================================
\node[anchor=north west, font=\sffamily\footnotesize\bfseries, text=blue!65!black]
  at (\xBL+0.20, \yC+\hC-0.05)
  {Banda de Control\enspace $f_{\mathrm{ctrl}}=50\,\mathrm{Hz}$,\enspace $T_c=0{,}02\,\mathrm{s}$};

\node[anchor=north east, font=\sffamily\footnotesize\bfseries, text=green!55!black]
  at (\xBR-0.20, \yF+\hF-0.14)
  {Banda de F\'{\i}sica\enspace $f_{\mathrm{phys}}=100\,\mathrm{Hz}$,\enspace $T_p=0{,}01\,\mathrm{s}$};

\node[anchor=north west, font=\sffamily\footnotesize\bfseries, text=orange!70!black]
  at (\xBL+0.20, \yT+\hT-0.14)
  {Banda de Telemetr\'{\i}a\enspace $f_{\mathrm{tel}}=10\,\mathrm{Hz}$,\enspace $T_t=0{,}1\,\mathrm{s}$};

%% ====================================================================
%%  BLOQUES -- BANDA DE CONTROL
%% ====================================================================
\node[widebox, fill=gray!7] (ref) at (-6.0, \yC+0.72) {
  \textbf{Referencia}\\
  $\bm{x}_{\mathrm{ref}},\,\dot{\bm{x}}_{\mathrm{ref}},\,\ddot{\bm{x}}_{\mathrm{ref}},\,\psi_{\mathrm{ref}}$
};

\node[box, fill=blue!7] (obs) at (-6.0, \yC-0.76) {
  \textbf{Observaci\'on}\\
  Ruido de medida\\
  $\bm{x}_{\mathrm{obs}},\,\dot{\bm{x}}_{\mathrm{obs}}$
};

\node[widebox, fill=blue!10] (ctrl) at (-1.0, \yC) {
  \textbf{Controlador}\\
  Cl\'asico o Neuronal\\
  $T,\;\bm{\tau}_B$
};

\node[box, fill=blue!7] (mix) at (5.2, \yC) {
  \textbf{Mezclador}\\
  Asignaci\'on y clipping\\
  $\bm{\omega}_{\mathrm{target}}$
};

%% -- Flechas control --
%%  ref -> ctrl: trayectoria deseada (salida horizontal, codo ascendente)
\draw[arr] (ref.east) -- ([xshift=0.5cm]ref.east) |- node[below, font=\sffamily\scriptsize, pos=0.75]{trayectoria deseada} (-1.0, 1.4) -- (ctrl.north);
%%  obs -> ctrl: estado observado (simplificado con ruteo absoluto robusto)
\draw[arr] (-4.4, -0.76) -- (-4.1, -0.76) |- node[above, font=\sffamily\scriptsize, pos=0.73]{estado observado} (-1.5, -1.58) -- (-1.5, -0.6);
%%  ctrl -> mix: comandos de empuje y momentos
\draw[arr] (ctrl.east) -- node[above, font=\sffamily\scriptsize]{comandos $T,\;\bm{\tau}_B$} (mix.west);

%% ====================================================================
%%  BLOQUES -- BANDA DE FISICA
%% ====================================================================
\node[box, fill=green!9] (act) at (5.2, \yF) {
  \textbf{Actuadores}\\
  Din\'amica de rotores\\
  retardo y lag
};

\node[widebox, fill=green!11] (rk4) at (-2.2, \yF) {
  \textbf{Integraci\'on RK4}\\
  Cuerpo r\'igido\\
  $\bm{x},\,\dot{\bm{x}},\,\bm{q},\,\bm{\omega}$
};

%%  Se ajusta pos=0.30 para posicionar la etiqueta en el espacio holgado a la derecha de la telemetria vertical (x=0.8)
\draw[arr] (act.west) -- node[above, font=\sffamily\scriptsize, pos=0.30]{velocidad real $\bm{\omega}_{\mathrm{rotor}}$} (rk4.east);

%% ====================================================================
%%  BLOQUE -- BANDA DE TELEMETRIA
%%  Centrado en x=2.0 para dejar espacio libre a la izquierda (ruta RK4->tel)
%% ====================================================================
\node[widebox, minimum width=5.2cm, fill=orange!6] (tel) at (2.0, \yT) {
  \textbf{Registro de Telemetr\'ia}\\
  Muestreo de estado y comandos
};

%% ====================================================================
%%  CANAL DERECHO (ZOH):  Mezclador -> Actuadores
%% ====================================================================
\draw[arr, draw=red!65!black, line width=1.2pt]
  (mix.east) -- (\xZOH, \yC)
  -- node[right, font=\sffamily\scriptsize\bfseries, text=red!75!black, align=left]
       {consigna $\bm{\omega}_{\mathrm{target}}$\\(ZOH 50\,Hz)}
     (\xZOH, \yF)
  -- (act.east);

%% ====================================================================
%%  CANAL IZQUIERDO (FB):  RK4 -> Observacion  (realimentación de estado real)
%%  xFB = -8.4, dentro de la banda de control, sube limpio por fuera de ref/obs
%% ====================================================================
\draw[darr]
  (rk4.west) -- (\xFB, \yF)
  -- node[right, font=\sffamily\scriptsize\itshape, align=left]
       {realimentaci\'on\\de estado real\\$\bm{x},\,\dot{\bm{x}},\,\bm{q},\,\bm{\omega}$}
     (\xFB, \yC-0.76)
  -- (obs.west);

%% ====================================================================
%%  TELEMETRIA 1:  RK4 -> Registro  (L limpia mediante ruteo ortogonal automático)
%%  Ruta: rk4.south (shift derecho) -> y=-7.9 -> tel.west
%% ====================================================================
\draw[tarr]
  ([xshift=1.0cm]rk4.south) |- node[left, font=\sffamily\scriptsize\itshape, pos=0.55, align=right] {registro de\\estado real}
  (tel.west);

%% ====================================================================
%%  TELEMETRIA 2:  Controlador -> Registro  (bajada por el canal central)
%%  Ruta: ctrl.south (shift) -> y=-2.35 -> x=0.8 -> y=-7.3 (tel.north)
%%  Cruzando la línea física (act.west -> rk4.east) en y=-4.6, x=0.8
%%  Para evitar solapamientos, dibujamos un arco de puente de salto en y=-4.6
%% ====================================================================
\draw[tarr]
  (-0.5, -0.6) -- (-0.5, -2.35)
  -- (0.8, -2.35)
  -- (0.8, -4.45)
  arc[start angle=90, end angle=-90, radius=0.15]
  -- (0.8, -7.3)
  node[right, font=\sffamily\scriptsize\itshape, pos=0.45, align=left] {registro de\\comandos};


\end{tikzpicture}
\end{document}
